\documentclass{article}
\usepackage{tikz}
\usetikzlibrary{shapes.geometric,  shapes.misc, arrows.meta, positioning}
\usepackage{graphicx}
\usepackage{url}
\usepackage{amsmath}
\usepackage{makeidx}
\usepackage{amssymb}

\begin{document}
\pagestyle{empty}


%\begin{figure}
%\vspace{-3cm}
%  \caption{Schema of the algorithm used in the model. Grey boxes are statistics obtained from the model in each step.}
%  \label{fig:diagram}
%\end{figure}

\begin{tikzpicture}[
remember picture, overlay,
  node distance=1cm and 1cm,
  >=Stealth,
  trap/.style={
    trapezium,
    draw,
    trapezium left angle=-70,
    trapezium right angle=-110,
    minimum width=4cm,
    minimum height=1.5cm,
    align=center,
    font=\small,
    inner sep=8pt
  },
  every node/.style={font=\small},
  block/.style={rectangle, draw, rounded corners, align=center, minimum width=3cm, minimum height=1cm},
  decision/.style={diamond, draw, aspect=2, align=center},
  terminator/.style={ellipse, draw, align=center
 },
 textnode/.style={rectangle, align=left, minimum width=2.5cm, minimum height=0.8cm, font=\footnotesize},
 statistics/.style={rectangle, fill=lightgray!30, align=center, minimum width=1.5cm, minimum height=0.4cm, font=\scriptsize\ttfamily},
 iterate/.style={isosceles triangle,
  draw,aspect=2,
  fill=gray!5,
  shape border rotate=0,
  minimum height=1.5cm,
  minimum width=2cm,
  inner sep=4pt,
  font=\small,
  align=center},
 decision/.style={diamond,
  draw,
  aspect=2.5, shape border rotate=0,
  minimum height=1cm,
  minimum width=1cm,
  inner sep=3pt,
  font=\small,
  align=center}
]
\begin{scope}[shift={(current page)}, xshift=-6cm, yshift=10.5cm]
\node[terminator] (start) {
  \texttt{model.initialize()} };
\node[iterate, below=of start] (sim) {
  \texttt{model.simulate\_full()} \\$\forall t=0..T$};
\node[block, below=of sim] (init) {
  \texttt{initialize\_step()}  };
\node[block, below=of init] (links) {
  \texttt{setup\_links()} };
\node[block, below=of links] (shock1) {
  \texttt{do\_shock("shock1")}  };
\node[block, below=of shock1] (loans) {
  \texttt{do\_loans()} };
\node[block, below=of loans] (shock2) {
  \texttt{do\_shock("shock2")} };
\node[block, below=of shock2] (repayments) {
  \texttt{do\_repayments()} };
\node[block, text=blue, below=of repayments] (ir) {
  \texttt{do\_interest\_rate()} };
\node[terminator, right=of sim] (finish) {
  \texttt{model.finish()} };

%
%\node[statistics, anchor=north east, at=(shock1.south)] (shock1_stats) {
%  potential\_lenders };
%
%\node[statistics, anchor=north east, at=(loans.south)] (loans_stats1) {
%  psi\_lenders psi\_loans \\
%  active\_lenders \\
%  active\_borrowers};
%
%\node[statistics, right=0cm of loans_stats1] (loans_stats2) {
%  leverage loans\\
%  equity\_lenders\\
%  systemic\_leverage
%};
%
%\node[statistics, anchor=north east, at=(repayments.south)] (repayments_stats1) {
%  liquidity best\_lender \\
%  bankruptcies \\
%  equity fitness};
%
%\node[statistics, right=0cm of repayments_stats1] (repayments_stats2) {
%  policy deposits \\
%  potential\_credit\_channels \\
%  reserves
%};
%\node[statistics, anchor=north east, at=(ir.south)] (ir_stats) {
%  P P\_max P\_min P\_std \\
%  prob\_bankruptcy \\
%  asset\_i asset\_j \\
%  interest\_rate };


%\node[textnode, right=of start] (start_text) {
%    Initial conditions for each bank $r^i_0=0.02$:
%    \begin{tabular}{l|l}
%        \multicolumn{1}{c|}{\textbf{Assets}} &  \multicolumn{1}{c}{\textbf{Liabilities}} \\ \hline
%        $C^i_0= 30-2.7 = 27.3$ &  $D^i_0= 135$\\
%        $L^i_0= 120$  & $E^i_0= 15$ \\
%        $R= 0.02 D^i_0 = 2.7$  &
%    \end{tabular}  };


%
%\node[textnode, right=of init] (init_text) {
%    $bank.B=0$ \\
%    $bank.rationing=0$ \\
%    $bank.active\_borrowers=0 ... $
%    };

\node[textnode, text=blue, right=of links] (links_text) {
    Assign new link depending on Boltzmann probability: $
      P_t^j = \frac{1}{1 + e^{-\beta(\mu_t^k - \mu_t^i)}}
    $ };
%
%\node[textnode, right=of shock1] (shock1_text) {
%  $ D^i_t = D^i_{t-1} \left( \mu + \omega U(0,1) \right )$
% \\
%  borrower if $\Delta D^i_t+C^i_t \le 0$ with demand $d=\Delta D$
%  \\
%  lender if $\Delta D^i_t+C^i_t > 0$ with supply $s= C +\Delta D$
% };
%
%\node[textnode, right=of shock1_text] (shock2_text) {
%    \begin{tabular}{l|l}
%        \multicolumn{1}{c|}{\textbf{Assets}} &  \multicolumn{1}{c}{\textbf{Liabilities}} \\ \hline
%        $C^i_t$ &  $D^i_t$\\
%        $L^i_t$  & $E^i_t$ \\
%        $R^i_t$  & $d$ (demand of loan)
%    \end{tabular}  };
%
%\node[decision, right=of loans] (loans_text1) {
%  $bank.d>0$? };
%\node[decision, right=of loans_text1] (loans_text2) {
%  $bank.d \le lender.s$? };
%\node[block, right=of loans_text2] (loans_text3) {
%  \texttt{do\_firesales()}};
%
%\node[block, below=0.3cm of loans_text2] (loans_text4) {
%  Obtain loan};
%
%\draw[->] (loans) -- (loans_text1);
%\draw[->] (loans_text1) -- node[above, font=\scriptsize] {Yes: it's borrower}  (loans_text2);
%\draw[->] (loans_text1.north) -- ++(0,0.4) node[above, font=\scriptsize]{No: it's lender} --   ++ (-3,0) --   (loans);
%\draw[->] (loans_text2) -- node[above, font=\scriptsize] {No}  (loans_text3);
%\draw[->] (loans_text2.south) -- node[above, font=\scriptsize] {Yes}  (loans_text4);
%\draw[->] (loans_text3) -- ++(0,-1.7) -- (loans_text4);
%
%\node[textnode, right=of shock2] (shock2_text) {
%  $ D^i_t = D^i_{t-1} \left( \mu + \omega U(0,1) \right )$
% };
%
%
%
%\node[decision, right=0.3cm of repayments] (repayments_text1) {
%  $bank.l>0$? };
%\node[decision, right=2cm of repayments_text1] (repayments_text2) {
%  $loan+interest<C$? };
%\node[block, right=1cm of repayments_text2] (repayments_text3) {
%  \texttt{do\_firesales()}};
%\node[decision, below=0.4cm of repayments_text2] (repayments_text4) {
%  $bank.d>0$?};
%
%\node[textnode, left=0.1cm of repayments_text4] (repayments_text7) { \\
%Equity updated: \\
%$E_t^i = E_{t-1}^i + \sum_{j} l_{t-1}^{i,j} r_{t-1}^{i,j} - \sum_{j \subseteq \theta_t^i} B_t^{i,j} - (1 - \rho) \tilde{L}_t^i$
%    };
%\node[decision, below=0.2cm of repayments_text4] (repayments_text5) {
%  $bank.failed$?};
%\node[block, right=0.4cm of repayments_text5] (repayments_text6) {
%  \texttt{replace\_bank()}};
%
%
%\draw[->] (repayments) -- (repayments_text1);
%\draw[->] (repayments_text1) -- node[above, font=\scriptsize] {Yes: pay loan}  (repayments_text2);
%\draw[->] (repayments_text1) -- node[above, font=\scriptsize] {No}  (repayments_text4);
%\draw[->] (repayments_text2) -- node[above, font=\scriptsize] {No}  (repayments_text3);
%\draw[->] (repayments_text4) -- node[below, font=\scriptsize] {Yes: shock2 negative and no capital}  (repayments_text3);
%\draw[->] (repayments_text4) -- node[below, font=\scriptsize] {No}  (repayments_text5);
%\draw[->] (repayments_text5) -- node[below, font=\scriptsize] {Yes}  (repayments_text6);
%

\node[textnode, right=0.3cm of ir] (ir_text1) {
  \quad \quad \quad \quad \quad \quad  $p_t^j=\frac{E^j_t}{E^{max}_t}$ \\
  $c^{i,j}_t=(1-h^i_j)A^j_t > 0$, if $(i,j) \in V_t$ \\
  $c^{i,j}_t=0$, otherwise \\
  $h^j_t=\frac{\lambda^j_t}{\lambda^{max}_t} $ \quad \quad \quad \quad
  $\lambda ^j_t=\frac{L^j_t}{E^j_t}$
 };
\node[textnode, right=0.3cm of ir_text1] (ir_text2) {
  Rate for next period of time: \\
  $r_t^{i,j} = \frac{\chi A_t^i - \phi A_t^j - (1 - p_t^j)(\xi A_t^j - c_t^{i,j})}{p_t^j c_t^{i,j}} (1 + \psi)$
 };
\node[textnode, text=blue, below=0.3cm of ir_text2] (ir_text3) {
  Profits: \\
  $\mu_t^i =  \eta_t \frac{C_t^i}{C^{max}_t}  + (1 - \eta_t) \frac{r^{min}_t}{r_t^i} $
 };
\draw[->] (ir) -- (ir_text1);
\draw[->, shorten >=1pt, shorten <=1pt] (ir_text1) -- (ir_text2);
\draw[->] (ir_text2) -- (ir_text3);





\draw[->] (start) -- (sim);
\draw[->] (sim) -- (init);
\draw[->] (init) -- (links);
\draw[->] (links) -- (shock1);
\draw[->] (shock1) -- (loans);
\draw[->] (loans) -- (shock2);
\draw[->] (shock2) -- (repayments);
\draw[->] (repayments) -- (ir);
\draw[->] (sim.east) -- (finish);
\draw[->] (ir.south) |- ++(0,-1.5) -- ++(-3,0) |- (sim.west);

\end{scope}
\end{tikzpicture}

\end{document}